\section{The sharp all-compact-test return criterion}
\label{sec:uniform-partial-return}

MA-PR-006 characterizes the return for one test by its
actual finite-block costs.  Uniform boundedness turns that
criterion into a precise statement for every compact test.
It also permits a single linear choice of the return maps.
The result below does not assert that its arithmetic cost
estimates hold, or that they fail.

Retain the fixed function, Hilbert parameters, actual
quotients, and block norms of MA-PR-006.
For a compact set \(K\subset\mathbb R\), put
\[
 \begin{aligned}
 \mathcal D_K
    &=\{g\in C^\infty(\mathbb R):\operatorname{supp}g\subset K\},\\
 p_j(g)&=\max_{0\le l\le j}\sup_{x\in K}|g^{(l)}(x)|.
 \end{aligned}
 \tag{MA-PR.81}
\]
Fix either of the following closed test spaces:
\[
 X_K=\mathcal D_K\cap\mathcal E_{\rm el},
 \qquad\hbox{or}\qquad
 X_K=\mathcal D_K\cap\mathcal E_{\rm el}
                              \cap\{g:g(0)=0\}.
 \tag{MA-PR.82}
\]
The choice may include the norm-one condition throughout.
All topological conclusions will be stated on these
fixed-support Fr\'echet spaces.
For \(g\in X_K\), define the possibly infinite seminorm
\[
              Q_k(g)=\sup_n\nu_n^k c_n(g),\qquad k\ge0,
 \tag{MA-PR.83}
\]
where \(c_n(g)\) is the attained minimum in (MA-PR.50),
measured on the whole \(Q_H\).

\paragraph{MA-PR-009 (uniform boundedness and a coherent linear return).}
For every fixed compact \(K\), the following are equivalent.
\begin{enumerate}
\item Every \(g\in X_K\) admits an actual continuous
equivariant return \(S_g:Q_H\to Q_\infty\) satisfying
(MA-PR.51).
\item For every integer \(k\ge0\), there are an integer
\(j=j(K,k)\) and a finite constant \(C=C(K,k)\) such that
\[
                   Q_k(g)\le C p_j(g)\qquad(g\in X_K).
 \tag{MA-PR.84}
\]
\item The returns can be chosen linearly in \(g\), with
actual linear lifts \(\widetilde S_g:Q_H\to\mathcal E\),
so that
\[
           (g,v)\longmapsto\widetilde S_gv
                     \quad\hbox{is jointly continuous on }
                     X_K\times Q_H .
 \tag{MA-PR.85}
\]
They satisfy the actual factorization and equivariance
of (MA-PR.51), and their common nuclear representations
have coefficients in every \(\ell^p\), \(p>0\), uniformly
when \(g\) ranges over a bounded subset of \(X_K\).
\end{enumerate}
The finite linear sections used for the third condition
may be chosen once for the fixed Hilbert parameters,
independently of \(K\).
Consequently, if these equivalent conditions hold on
every compact-support space, there is one linear return
assignment on all eligible compact tests in the chosen class, with the
joint continuity and nuclear bounds just stated on
each \(X_K\).
It suffices to test the intervals \(K=[-R,R]\),
\(R=1,2,\ldots\).

There is a sharp failure alternative.  Suppose that
for one fixed \(K\) and one fixed \(k\), no estimate
(MA-PR.84) holds for any finite \(j,C\).
Then
\[
       \{g\in X_K:Q_k(g)=\infty\}
          \quad\hbox{is a dense }G_\delta\hbox{ subset of }X_K .
 \tag{MA-PR.86}
\]
Every test in that set has no continuous equivariant
return.  This is a conditional topological alternative,
not a claim that the displayed failure premise occurs
for an arithmetic \(L\)-function.

\begin{proof}
\emph{The exact Fr\'echet space and its cost seminorms.}
The increasing seminorms in (MA-PR.81) make
\(\mathcal D_K\) a complete metrizable locally convex space.
Indeed uniform convergence of every derivative on a
fixed interval containing \(K\) gives a smooth limit,
by integrating the derivative limits; the limit vanishes
off \(K\).  All endpoint flatness conditions are retained.
Each Mellin jet in (MA-PR.49) is continuous on this
space, because its integral is over the fixed compact
support.  Evaluation at zero is continuous as well.
Thus either space in (MA-PR.82) is a closed Fr\'echet
subspace.

On \(X_K\), every \(c_n\) is a continuous seminorm.
Homogeneity follows by scaling feasible module maps.
For the triangle inequality, the sum of a feasible
map for \(g\) and one for \(h\) is feasible for \(g+h\),
so minimizing the triangle inequality for their actual
operator norms gives
\[
        c_n(g+h)\le c_n(g)+c_n(h),\qquad
        c_n(tg)=|t|c_n(g).
 \tag{MA-PR.87}
\]
For continuity, the fixed linear selector of
(MA-PR.63)--(MA-PR.64) gives
\(c_n(g)\le A_n\|b_n(g)\|_n\).
If \(K\subset[-R,R]\), the fixed full strip of the
M18 block contours gives
\[
 \|b_n(g)\|_n
       \le \sup_{w\in\partial\Omega_n}
                   \left|\int_K e^{(w-a)x}g(x)\,dx\right|
       \le 2R e^{(2+|a|)R}p_0(g).
 \tag{MA-PR.88}
\]
Take a positive \(R\) even for a degenerate compact set.
For each fixed \(n\), \(A_n\) is finite; no uniform
growth bound for these constants is used.
This proves continuity of \(c_n\).
The supremum \(Q_k\) is therefore a lower
semicontinuous extended seminorm.

\emph{Pointwise returns imply finite-seminorm bounds.}
Assume the first condition.
MA-PR-006 makes \(Q_k(g)<\infty\) for every \(g\in X_K\)
and every fixed \(k\).
For this \(k\), let
\[
             E_M=\{g\in X_K:Q_k(g)\le M\},
                         \qquad M=1,2,\ldots .
 \tag{MA-PR.89}
\]
These are closed, balanced and convex, and they cover \(X_K\).
The Baire property of the complete metrizable space
makes one \(E_M\) have nonempty interior.

For clarity, the complete-metric fact used here can
be seen directly.  If countably many closed sets all
had empty interior, start inside an arbitrary nonempty
open set and choose successively nonempty closed balls,
each contained in the interior of its predecessor,
avoiding the next closed set, with radii tending to zero.
Their centers form a Cauchy sequence and their common
limit, by completeness, avoids every closed set.
Thus such sets cannot cover the space.
The same construction starting in any open set proves
the density of the complement of their union.

Choose \(g_0\) in the interior of this \(E_M\).
For some neighborhood \(U\) of zero,
\(g_0+U\subset E_M\).  Since \(g_0\in E_M\),
(MA-PR.87) gives \(Q_k(u)\le2M\) for \(u\in U\).
The increasing Fr\'echet seminorms imply that \(U\)
contains \(\{g:p_j(g)<\epsilon\}\) for some \(j\)
and \(\epsilon>0\).
Scaling, and then letting the strict inequality
approach \(\epsilon\), gives
\[
                  Q_k(g)\le (2M/\epsilon)p_j(g).
 \tag{MA-PR.90}
\]
This is (MA-PR.84).  The index \(j\) and the constant
may depend on \(K\) and \(k\).

\emph{A finite linear section with the precise dimension loss.}
We next construct the coherent family required for
the third condition.  This avoids selecting minimizers
nonlinearly as the test varies.
Define the finite-dimensional normed space
\[
 \mathcal X_n=\operatorname{Hom}_{\mathbb C[w]}(B_n,J_n),
 \qquad
       \|L\|_{\mathcal X_n}=\|L\mathscr R_n\| .
 \tag{MA-PR.91}
\]
This is a genuine norm since \(\mathscr R_n\) is
surjective.  All source norms remain those on the
whole Hilbert quotient.
Let
\[
 \begin{aligned}
 \Delta_n:\mathcal X_n&\longrightarrow
                 \operatorname{End}_{\mathbb C[w]}(B_n),
             &\Delta_nL&=\pi_nL,\\
 \mathcal Y_n&=\operatorname{im}\Delta_n,
             &\|y\|_{\mathcal Y_n}
                 &=\min_{\Delta_nL=y}\|L\|_{\mathcal X_n}.
 \end{aligned}
 \tag{MA-PR.92}
\]
This is the actual quotient norm, with attained
minimum by the proof of (MA-PR.53).
An endomorphism of \(B_n\) is multiplication by
its value on the cyclic vector \(1\).
The primary calculation (MA-PR.54) therefore gives
\[
 e_n:=\dim\mathcal Y_n
       =\sum_{\rho\in\Omega_n}
                     \max(0,2r_\rho-m_\rho)
       \le\dim B_n\le d_n=O(\log\nu_n).
 \tag{MA-PR.93}
\]
At a full retained factor the summand is \(m_\rho\);
at a discarded factor it is zero.

Choose an Auerbach basis \(y_{n,l}\), \(y^*_{n,l}\)
of \(\mathcal Y_n\): all vectors and functionals
have norm one.  The elementary determinant-maximization
construction of MT-M18 supplies this basis.
For each \(y_{n,l}\), choose an attained minimum
preimage \(L_{n,l}\), of \(\mathcal X_n\)-norm one.
Then
\[
 \begin{aligned}
 \sigma_n(y)&=\sum_{l=1}^{e_n}y^*_{n,l}(y)L_{n,l},&
 \Delta_n\sigma_n&=1,\\
 \|\sigma_n(y)\|_{\mathcal X_n}
                  &\le e_n\|y\|_{\mathcal Y_n}.
 \end{aligned}
 \tag{MA-PR.94}
\]
If \(e_n=0\), use the zero map.
These are fixed complex-linear sections of the
finite spaces of feasible homomorphism data.
They are not module sections of
\(\pi_n:J_n\to B_n\), which a proper truncation
would forbid.

For an eligible test, set
\[
 \begin{aligned}
 y_n(g)&=M_{\pi_nb_n(g)}\in\mathcal Y_n,&
 \|y_n(g)\|_{\mathcal Y_n}&=c_n(g),\\
 \widehat F_{n,g}&=\sigma_n(y_n(g))\mathscr R_n .
 \end{aligned}
 \tag{MA-PR.95}
\]
These maps are linear in \(g\) and satisfy
\[
 \begin{aligned}
 \widehat F_{n,g}&=\widehat F_{n,g}P_n,&
 \|\widehat F_{n,g}\|&\le e_nc_n(g),\\
 \pi_n\sigma_n(y_n(g))&=M_{\pi_nb_n(g)} .
 \end{aligned}
 \tag{MA-PR.96}
\]
Thus the dimension loss is at most \(O(\log\nu_n)\),
with no hidden interpolation or edge-distance constant.

\emph{Uniform estimates give the actual continuous family.}
Assume (MA-PR.84) for every \(k\).
Define
\[
        \widetilde S_gv
              =\sum_n\mathcal T_n\widehat F_{n,g}v,
              \qquad S_g=q_\infty\widetilde S_g .
 \tag{MA-PR.97}
\]
For every defining seminorm \(p\) of \(\mathcal E\),
choose a sufficiently large integer \(k_p\) relative
to its fixed polynomial lift exponent in (MA-PR.44).
Equations (MA-PR.93) and (MA-PR.96) give
\[
 \begin{aligned}
 p(\widetilde S_gv)
   &\le C_p\sum_n\nu_n^{D_p}e_nc_n(g)\|v\|\\
   &\le C'_p Q_{k_p}(g)\|v\|
    \le C''_{K,p}p_{j(K,k_p)}(g)\|v\|.
 \end{aligned}
 \tag{MA-PR.98}
\]
For example any integer \(k_p>D_p+2\) makes the
sum converge.  The series is absolutely convergent
in every seminorm, uniformly on the indicated
bounded sets.  This proves its linearity in \(g\)
and the joint continuity in (MA-PR.85).
Its maps satisfy the actual factorization and
equivariance by the same full-jet and actual
common-kernel argument as (MA-PR.58).

Let \(\mathcal B\subset X_K\) be bounded.
The estimates (MA-PR.84) for every \(k\) show that
\[
       \beta_n=e^{-\nu_n}
                   +\sup_{g\in\mathcal B}\|\widehat F_{n,g}\|
                   \quad\hbox{decreases rapidly}.
 \tag{MA-PR.99}
\]
For an Auerbach basis \(e_{n,j},\ell_{n,j}\) of \(J_n\),
use target vectors
\(\beta_n^{1/2}\mathcal T_ne_{n,j}\),
coefficients \(\beta_n^{1/2}\), and source functionals
\(\ell_{n,j}\widehat F_{n,g}/\beta_n\).
The source norms are at most one for all
\(g\in\mathcal B\), the target vectors form one
bounded set in \(\mathcal E\), and the coefficients,
repeated \(d_n\) times, belong to every \(\ell^p\).
This proves exactly the asserted uniform
all-order nuclear representation.
There is no passage from pointwise bounds to
an unjustified uniform near-boundary estimate:
the only uniform step is the proved Baire implication
(MA-PR.90).

The sections \(\sigma_n\) depend only on the fixed
finite normed spaces, so the same formula (MA-PR.97)
works on every compact-support space for which
(MA-PR.84) holds.  This proves coherence across
supports.  Continuity is asserted on each \(X_K\);
no additional assertion about a different inherited
LF topology is needed.
The third condition plainly implies the first,
which completes the equivalence.

\emph{The residual failure alternative.}
Fix \(K,k\) and suppose that every estimate in
(MA-PR.84) fails.
The closed sets \(E_M\) of (MA-PR.89) then have
empty interior: otherwise the exact neighborhood
and scaling argument proving (MA-PR.90) would
give one of the forbidden finite-seminorm estimates.
Their open complements are dense.
The complete-metric argument given above shows
that
\[
       X_K\setminus\bigcup_{M\ge1}E_M
          =\bigcap_{M\ge1}(X_K\setminus E_M)
\]
is a dense \(G_\delta\).  This set is precisely
\(\{g:Q_k(g)=\infty\}\).
MA-PR-006 excludes a continuous equivariant return
for every one of its elements.  This proves
(MA-PR.86) without constructing or assuming
an arithmetic close-pair sequence.
\end{proof}

\paragraph{Exact remaining gate and source placement.}
The all-compact-test question is now exactly the
family of finite derivative estimates (MA-PR.84).
The theorem proves neither these estimates nor
their failure for the arithmetic function.
MA-PR-008 supplies one possible necessary quantitative
test of them, conditional on an actual retained/discarded
configuration; it supplies no such configuration.
The adaptive returns of MA-PR-007 remain available
for every eligible test without (MA-PR.84).

The spaces \(J_n\) and \(B_n\) remain actual
finite-length \(\mathbb C[w]\)-modules.
The spaces \(\mathcal X_n\) and \(\mathcal Y_n\)
in this argument organize feasible module maps
and carry their measured operator and quotient norms.
The Auerbach section (MA-PR.94) is only complex linear
in this space of maps.  It does not split the local
module extension (MA-PR.80), and no splitting is
inferred from its relation in the Grothendieck group
of finite-length modules.
The common nuclear series in (MA-PR.97)--(MA-PR.99)
is the actual analytic realization; finite traces
and exterior-power determinants retain the scope
already specified in MA-PR-006--008.
No new arithmetic positivity or zero-location claim
is obtained.

The exact inputs are MA-PR-006, including its
whole-quotient Gram norm, minimum attainment and
actual return criterion; the finite module calculation
MA-PR.54; the full smooth adapted topology
MT-M18 (M18.1)--(M18.3), its polynomial entire-lift
bound (M18.13), and its complete Auerbach-basis
construction.  The Baire implication and the
finite linear-section estimate used here are proved
above.  The inherited cote~1 nuclear-trace locators
remain scans~079--080 and~111--112, as in MA-WP
and MA-PR-006--008; no new reading of those scans
is claimed.  No A, M, generic-edge assumption,
or new outside preprint is an input.
